
\section{Discussion: why the detectors disagree}
\label{sec:discussion}

Across three enacted congressional plans our learned geometric detector and
the classical vote-asymmetry statistics disagree in a patterned way, and the
pattern has a single mechanistic source. A partisan advantage can be built two
ways. Cracking distributes an opponent's voters across many districts by
threading boundaries through scattered pockets of support, and packing
concentrates them into a few districts, often by exploiting voters who already
live in compact clusters. Cracking distorts district shape and severs
adjacency, whereas packing a naturally clustered population requires little or
no boundary distortion. A detector that reads only geometry is therefore
sensitive to the first mechanism and blind to the second, and this asymmetry
organizes every result below.

\paragraph{Geometry detects cracking and is blind to packing.}
The topology-only model flags Pennsylvania's 2011 plan ($p=0.009$) and
Maryland's 2011 plan ($p=0.016$) but not North Carolina's ($p=0.36$), and a
direct compactness audit against each state's neutral ensemble explains why.
Measured by cut edges, the number of dual-graph edges severed by district
boundaries and the quantity a message-passing network is most directly
sensitive to, the Pennsylvania plan sits at the ensemble maximum ($+5.8$
standard deviations) and the Maryland plan carries more than double the cut
edges of any neutral map ($+13$ standard deviations), while North Carolina's
plan is entirely typical (the $42.6$th percentile) and is in fact more compact
than the neutral ensemble on the Polsby--Popper score. Cut edges and
contour-based compactness are properties of how a boundary encloses area,
independent of which voters live inside it \cite{duchin2023discrete,
validi2022cutedges}, and a compact map can nonetheless be a severe partisan
gerrymander: Alexeev and Mixon prove a party can win supermajorities of seats
while every district remains geometrically compact \cite{alexeev2018partisan},
and compactness scores are unstable to fairness-irrelevant implementation
choices and can be gamed so that gerrymandered districts appear reasonable
\cite{barnes2021gerrymandering, polsby1991third}. North Carolina's advantage is
obtained by packing opposition voters into already-compact urban clusters, so
its plan leaves no geometric signature and the shape detector correctly reports
nothing; Pennsylvania's and Maryland's advantages depend on boundary
distortion, which places them in the tail of the geometric null. The detector
is thus complementary to, not a replacement for, vote-based statistics: it sees
the cracking mechanism they can miss and is silent on the packing mechanism
only they can see.

\paragraph{Vote-asymmetry statistics fail under a lopsided baseline.}
The mean--median difference flags North Carolina under all five available
statewide elections, Pennsylvania under four of nine, and Maryland under none
of six, and the failure tracks each state's statewide partisan balance
(two-party Democratic shares of $0.465$, $0.527$, and $0.609$). The statistic
is by construction insensitive to any reallocation of votes among same-side
districts that leaves the median district unchanged \cite{warrington2018quantifying},
which is precisely Maryland's configuration: a wall of safe but not
overwhelming Democratic seats, the structure that vote-asymmetry measures are
known to miss \cite{nagle2020unbalanced}. This is not incidental to the
statistic but intrinsic to it. Its interpretation as a bias signal is inherited
from a uniform-partisan-swing counterfactual \cite{katz2020theoretical}, and
DeFord and Veomett prove that the achievable range of mean--median and related
symmetry scores collapses as the statewide vote moves away from parity, so
under a sufficiently lopsided baseline the statistic cannot register an extreme
seat outcome regardless of how gerrymandered the map is \cite{deford2024bounds}.
An ensemble study using the same recombination sampler we adopt finds the same
degradation empirically, with mean--median distinguishing seat extremes far
better in a balanced state than a lopsided one \cite{deford2021implementing,
deford2021recombination}. Declination, which normalizes each party's average
winning margin by its own seat share rather than comparing raw shares to a fixed
threshold, is more robust to this failure \cite{warrington2018comparison} and
in our sweep recovers Maryland under four of six elections where mean--median
recovers none.

\paragraph{Confounded features act as a shortcut and degrade the model.}
Supplying the network with density, minority voting-age-population share, and
Democratic vote share as node features makes it worse, not better. The
augmented variant becomes $15$ to $72$ times more sensitive to a vote-share
label shuffle than the exactly invariant topology-only variant, and its
node-level localization collapses below chance on Maryland ($\mathrm{AUC}=0.49$)
where the topology-only variant succeeds ($\mathrm{AUC}=0.62$ to $0.74$).
Simplicity-biased optimization preferentially fits the simplest predictive
feature and remains invariant to more complex but equally predictive signal
\cite{shah2020simplicity, geirhos2020shortcut}, so the model latches onto the
vote-share correlate rather than integrating the structural signal; because that
correlate is spurious with respect to the redistricting mechanism, the reliance
sharpens with the feature's signal strength \cite{sagawa2020overparameterization}
and induces the out-of-distribution collapse observed when the model is asked to
generalize to a new state's geometry \cite{sagawa2020distributionally,
nagarajan2021failure}. The demographic augmentation is the settlement-density
confound entering the model directly, and the topology-only design is what keeps
the detector reading geometry rather than shortcutting through it.

\paragraph{A pooled score dilutes a localized anomaly.}
The plan-level anomaly score misses North Carolina ($p=0.36$) even though the
node-level readout on the same plan recovers its packed and cracked precincts.
Graph-level and node-level anomaly detection are distinct tasks precisely
because a single pooled statistic cannot preserve an anomaly confined to a small
subset of nodes \cite{ma2023survey}, which is the reason the distillation
objective we build on carries both a node-level and a graph-level term
\cite{ma2021glocalkd}. Aggregating per-node discrepancy into one plan score
averages a handful of packed precincts together with hundreds of ordinary ones,
driving the anomalous contribution toward the mean by the same
convergence-to-a-shared-value dynamic formalized for oversmoothing
\cite{rusch2023oversmoothing}; more generally, aggregate metrics are known to
hide failures visible only within a specific coherent subset
\cite{salaudeen2025aggregation}. The node-level readout acts before this
dilution, which is why localization succeeds where the pooled verdict does not.

\paragraph{Implications and limitations.}
Taken together the results position the learned geometric score as a third
family of ensemble-position tests, complementary to the vote-asymmetry
statistics and to hand-built geometric distances, that is sensitive to the
cracking-and-distortion mechanism and stable across elections because it never
consumes vote data. Its complementarity is exactly its limitation: it misses
North Carolina's packing-based gerrymander, and no single detector in our study
flags all three plans, which is itself an argument against reliance on any one
metric \cite{warrington2018comparison, deford2024bounds} and for reporting a
panel of geometric and vote-based ensemble tests together. Two points of legal
provenance deserve care. The Pennsylvania outlier analysis submitted in
\emph{League of Women Voters v.\ Commonwealth} concerned a remedial proposal
rather than the 2011 enacted plan our data encode, so we cite it for methodology
and for Pennsylvania's litigation history rather than as evidence about our
specific map. The Maryland plan we analyze is the 2011 congressional map whose
sixth district was the partisan-gerrymander claim in \emph{Lamone v.\ Benisek},
decided in \emph{Rucho v.\ Common Cause}; the third district, the ``praying
mantis,'' is geometrically notorious but was not itself adjudicated on shape
grounds, and we treat it accordingly.
